% !TEX root = ../main.tex
\chapter{Conclusion}
\label{ch:conclusion}
\label{sec:conclusion}

The book makes two connected claims and draws one practical conclusion.

The first claim is a redefinition. Hallucination, as ordinarily defined, is a false
or fabricated output, and that definition is wrong in both directions. It convicts a
system that faithfully reports an applicable source which happens to be
wrong---an error, whose fault lies in the evidence---and it acquits a system that
asserts something true while resting on nothing, which is the same defect as
fabrication with a luckier outcome. Defined instead as ungrounded assertion, with no
truth condition, hallucination becomes a property of the relation between an
assertion and an evidential base. That is not a terminological preference. The
truth-conditional definition can only be evaluated against reference answers, offline,
on questions already answered; the redefinition can be evaluated from the system's
own execution trace, at the point of assertion, and therefore used as a gate.

From the redefinition, what a system can assert falls into five categories, not two. Crossing
grounding against truth separates a right answer reached for a reason from a right
answer reached for none, and a faithful report of a bad source from a fabrication.
Requiring grounding to be applicable at a resolved context $c = \{j, t_v, f, u\}$
separates a third case from both: an assertion resting on evidence that is
authentic, versioned, and authoritative, and that does not govern the decision at
hand. These five categories reduce to three remedies with three different owners,
and a pooled hallucination rate identifies none of them.

The failures are structural. The architecture maps token sequences to
distributions over token sequences, with no stage at which two encodings of one
claim are brought to a common form. The pretraining objective contains no term that
is a function of truth, evidence, provenance, or context, and none minimized by
declining to answer. In-context learning makes the prompt the entire control
surface, and the prompt is a string. Preference optimization fits a reward to
comparisons between outputs, so grounding is not visible to it, and the pressure
runs toward assertion and away from abstention. And there is a statistical floor: a
calibrated model must assert some rare facts incorrectly, and enterprise questions
live where facts are rare.

The second claim---and it is the step that decides what to measure---is that there
is no canonical form. Nothing constructs the map to be constant on classes of materially equivalent
requests, and nothing in the objective penalizes divergence between them. So a
passing test case establishes that one string produced an acceptable output on one
occasion, and not that a materially equivalent string will, nor that the same string
will later in the same session. The equivalence classes are unbounded, so coverage
is not a remedy; the failures are silent, so nothing alerts; and repeatability fails
too, so the guarantee does not securely cover even the tested case. Semantic
normalization is not the fix, because material equivalence is context-sensitive,
depends on domain ontology, and turns on distinctions that regulated language uses
deliberately. Accuracy is a property of a sample, and the sample does not generalize
in the direction validation needs of it. A model's latent representation is
therefore not a candidate system of record: it has no stable identity, no
inspectable normalized form, no declared dependencies, and no reproducible
derivation.

Grounding does not have this defect. Checking correctness requires knowing
the answer, so it can happen only offline, on questions nobody needed to ask.
Checking grounding requires knowing what evidence the system used and whether it
governs the case---both properties of the system's own execution, available at
runtime, per assertion, on live traffic. That is the difference between a metric and
a gate, and blocking, escalating, qualifying, and logging a reconstructable basis all
require a gate. Proposition~\ref{prop:independence} shows that neither figure identifies the other: no accuracy figure constrains the rate at which a
system asserts beyond its evidence, and an intervention that raises accuracy by
encouraging confident guessing lowers that rate while improving the headline number.

The practical conclusion is a requirement on architecture rather than a
recommendation about metrics. The grounding rate is well defined for any system and
measurable only where the evidential base is explicit at inference
time---identified, versioned, scope-filtered, and resolved to a context, at claim
granularity, with verification of the support relation and an abstention path that
is a branch rather than a sentence. The nineteen invariants of
Section~\ref{sec:invariants} separate what the model will not supply from what the
system must, and the division of that work across data, policy, and platform
functions is itself a principal cause of failure: the operational unit of failure is
the ungoverned handoff, and the grounding rate is the right cross-functional metric
precisely because no single function can move it alone.

None of this is an argument against using these systems. I build them, and they
work. It is an argument about what to measure, in what order, and what a passing
test entitles an organization to conclude. A language model may be useful, fluent,
and often accurate. But in a regulated enterprise, a proposition becomes fit for
consequential reliance only when its context is resolved, its evidence is authentic
and applicable, its support relation is verified, its production path is
reconstructable, and its use is authorized within an accountable institutional
process.
